\documentclass[10pt,letterpaper]{article}

\usepackage[margin=0.60in]{geometry}
\usepackage[T1]{fontenc}
\usepackage{lmodern}
\usepackage{microtype}
\usepackage{enumitem}
\usepackage{titlesec}
\usepackage[hidelinks]{hyperref}
\usepackage{xcolor}
\usepackage{tabularx}
\usepackage{array}

% Content follows cv.html; keep both sources in sync when updating this CV.
\definecolor{ink}{RGB}{37,36,32}
\definecolor{muted}{RGB}{98,94,88}
\definecolor{rulegray}{RGB}{206,200,190}
\renewcommand{\familydefault}{\sfdefault}
\pagestyle{empty}
\setlength{\parindent}{0pt}
\setlength{\parskip}{0pt}
\setlength{\tabcolsep}{0pt}
\setlength{\emergencystretch}{1.5em}
\setlist[itemize]{leftmargin=1.25em,label=\textbullet,itemsep=2.5pt,topsep=4pt,parsep=0pt,partopsep=0pt}
\titleformat{\section}
  {\fontsize{9.5}{11.5}\selectfont\bfseries\color{ink}}
  {}{0pt}{\MakeUppercase}
  [\vspace{3pt}{\color{rulegray}\titlerule[0.4pt]}]
\titlespacing*{\section}{0pt}{13pt}{7pt}
\hypersetup{
  pdftitle={Curriculum Vitae -- Ziyu Peng},
  pdfauthor={Ziyu Peng},
  pdfsubject={Undergraduate research curriculum vitae}
}

\newcommand{\entry}[2]{%
  \begin{tabularx}{\textwidth}{@{}>{\raggedright\arraybackslash}X@{\hspace{12pt}}r@{}}
    \textbf{#1} & {\fontsize{9.5}{11.5}\selectfont\color{muted}#2}
  \end{tabularx}}
\newcommand{\detail}[1]{%
  {\fontsize{9.5}{11.8}\selectfont\color{muted}#1\par}}
\newcommand{\projectgap}{\vspace{6pt}}
\newcommand{\cvlink}[2]{\href{#1}{\underline{#2}}}

\begin{document}
\color{ink}
\fontsize{10}{12.2}\selectfont

\begin{tabularx}{\textwidth}{@{}Xr@{}}
  {\rmfamily\fontsize{29}{32}\selectfont\bfseries Ziyu Peng}
  & {\fontsize{9.5}{11.5}\selectfont\color{muted}Boulder, Colorado}
\end{tabularx}
\vspace{5pt}

{\fontsize{9.5}{11.5}\selectfont
Undergraduate Student \textperiodcentered{} Statistics \& Data Science \textperiodcentered{} University of Colorado Boulder\par}
\vspace{5pt}
{\fontsize{9.5}{11.5}\selectfont
\cvlink{mailto:Ziyu.PengSr@colorado.edu}{Ziyu.PengSr@colorado.edu}
\hspace{14pt}
\cvlink{https://pzy54154631-hub.github.io/homepage/}{Academic homepage}
\hspace{14pt}
\cvlink{https://github.com/pzy54154631-hub}{GitHub}\par}

\section*{Education}
\entry{University of Colorado Boulder}{2026--Present}
\detail{B.S. in Statistics and Data Science, in progress}
\vspace{6pt}
\entry{Xi'an Jiaotong-Liverpool University}{2025--2026}
\detail{First-year study in Data Science and Big Data Technology \textperiodcentered{} \textbf{GPA 3.88/4.00}}

\section*{Research \& Technical Projects}
\entry{Equity Return Prediction and Low-Turnover Stock Selection}{2026}
\detail{2026 Shenzhen International FinTech Competition (FinTechathon) \textperiodcentered{} Data Analytics Track}
\begin{itemize}
  \item Analyzed 7.9 million training observations from 4,650 A-share stocks; built 236 price-volume features and evaluated XGBoost and GRU models for cross-sectional return prediction.
  \item Developed separate stock-selection and daily-ranking layers. In 2022--2024 \emph{development} validation, reduced Top-group Jaccard turnover from 0.76--0.80 to 0.02--0.04, with a trade-off in estimated excess return.
\end{itemize}

\projectgap
\entry{Ambiguity Is Not One Scalar}{2026}
\begin{itemize}
  \item Implemented a visual question-answering router that enumerates plausible referents before choosing answer, clarify, or abstain.
  \item Evaluated a 26,425-example diagnostic benchmark; referent enumeration recovered 20 of 22 human-audited ambiguity misses versus 0 of 22 for an answer-consistency baseline. \cvlink{https://github.com/pzy54154631-hub/ambiguity-is-not-one-scalar}{Code and results}.
\end{itemize}

\projectgap
\entry{PW-OFPC: Object-Centric Predictive Coding}{2025--2026}
\begin{itemize}
  \item Developed an abstract visual-reasoning pipeline that retains uncertainty over object attributes and compares candidates using precision-weighted residuals.
  \item Reported RAVEN macro accuracy of 0.857 for cross-option inference versus 0.537 for direct prompting; maintained scripts to reproduce archived evaluation tables. \cvlink{https://github.com/pzy54154631-hub/pw-ofpc}{Code and results}.
\end{itemize}

\projectgap
\entry{Corridor Wealth: Auditable AI Advisory Prototype}{2026}
\begin{itemize}
  \item Built a synthetic-data financial advisory prototype with deterministic policy checks separate from model proposals and individually authorized simulated transactions. \cvlink{https://github.com/pzy54154631-hub/Shenzhen-Fintech-international-Project}{Project repository}.
\end{itemize}

\section*{Honors \& Recognition}
\textbf{FinTechathon 2026} --- Data Analytics Track: Top 25 (top 10\%), semifinal in progress; separately participating in the International AI Track.\par
\vspace{4pt}
\textbf{CU Boulder Transfer Excellence Scholarship} --- \$12,000 total (2026).\par
\vspace{4pt}
\textbf{XJTLU Academic Scholarship} (2026/27); \textbf{Third Prize}, XJTLU Biostatistics Competition (2026).\par

\section*{Technical Skills}
\textbf{Programming and tools:} Python, R, PyTorch, XGBoost, Git/GitHub, Linux, LaTeX.\par
\vspace{4pt}
\textbf{Methods:} Feature engineering, statistical modeling, temporal validation, ablation studies, and reproducible evaluation.

\end{document}
